\section{CLI Arguments Module}

\subsection{Purpose}

The args module is the single point of CLI parsing.
It decouples \texttt{main.cpp} from raw \texttt{argc}/\texttt{argv} handling and makes the
parsed configuration available to the rest of the program through a plain data struct.
All validation that can be done at the CLI level (e.g., checking that required flags are
present) happens here, so the codec functions can assume a coherent input.

\subsection{Files}

\begin{itemize}
  \item \texttt{include/args.hpp} --- \texttt{ParsedArgs} struct and \texttt{parse\_arguments}
        declaration.
  \item \texttt{src/args.cpp} --- implementation using \texttt{argparse.hpp}.
  \item \texttt{include/argparse.hpp} --- third-party header-only argument parsing library
        (bundled in-tree; see Section~\ref{sec:argparse}).
\end{itemize}

\subsection{ParsedArgs Struct}

\begin{lstlisting}[style=cpp]
struct ParsedArgs {
    bool        decompress    = false;
    std::string infile;
    std::string outfile;

    // Compression-only parameters.
    // During decompression these are IGNORED -- loaded from the file header.
    bool use_model    = false;
    bool adaptive_scan = false;
    int  width        = -1;
};
\end{lstlisting}

Field documentation:

\begin{center}
\begin{tabular}{>{\ttfamily}l l}
\toprule
\textbf{Field} & \textbf{Meaning} \\
\midrule
decompress     & \texttt{true} when \texttt{-d} flag is present; \texttt{false} for \texttt{-c} \\
infile         & Path to input file (from \texttt{-i}) \\
outfile        & Path to output file (from \texttt{-o}) \\
use\_model     & \texttt{true} when \texttt{-m} is present (compression only) \\
adaptive\_scan & \texttt{true} when \texttt{-a} is present (compression only) \\
width          & Image width in pixels from \texttt{-w}; default \texttt{-1} (invalid sentinel) \\
\bottomrule
\end{tabular}
\end{center}

\textbf{Important:} \texttt{use\_model}, \texttt{adaptive\_scan}, and \texttt{width} are only
meaningful during compression.
During decompression the codec ignores them entirely and reads the equivalent information
from the compressed file header (flags byte and width/height fields).
This design ensures the decompressor is self-contained.

\subsection{parse\_arguments()}

\texttt{ParsedArgs parse\_arguments(int argc, char** argv)} is the only exported function.
It performs the following steps:

\begin{enumerate}
  \item Registers all flags and options with \texttt{argparse.hpp}.
  \item Calls the library's \texttt{parse} method against \texttt{argc}/\texttt{argv}.
  \item Extracts values into a \texttt{ParsedArgs} instance.
  \item Validates mode consistency:
    \begin{itemize}
      \item Exactly one of \texttt{-c} or \texttt{-d} must be present.
      \item \texttt{-i} and \texttt{-o} are required in both modes.
      \item In compression mode (\texttt{-c}), \texttt{-w} must be provided and positive.
    \end{itemize}
  \item On any validation failure: prints a usage message to \texttt{stderr} and calls
        \texttt{exit(1)}.
\end{enumerate}

\subsection{argparse.hpp}
\label{sec:argparse}

\texttt{argparse.hpp} is a third-party header-only C++ argument parsing library bundled
directly in \texttt{include/}.
It handles option registration, type conversion, and help text generation.

It is \textbf{not} a compression library and does not violate the assignment's ``no external
compression libraries'' constraint.
Using it avoids reimplementing \texttt{getopt}-style parsing and reduces
\texttt{args.cpp} to \textasciitilde{}30 lines of straightforward setup code.
